\documentclass[11pt]{article}
\usepackage[utf8]{inputenc}
\usepackage{amsmath,amssymb}
\usepackage{hyperref}
\title{Robust Epigenetic Dna Methylation under Distribution Shift}
\author{Anonymous}
\date{}
\begin{document}
\maketitle

\begin{abstract}
This paper studies Epigenetic Dna Methylation. Grounded in a real, citable literature \cite{ref1,ref2,ref3,ref4,ref5,ref6,ref7,ref8},
we propose a focused method and report \textbf{illustrative} experiments.
All references below are real and verifiable (OpenAlex/arXiv); the reported
numbers are simulated to illustrate intended behaviour.
\end{abstract}

\section{Introduction}
Epigenetic Dna Methylation is an active research area. This work targets a concrete gap identified from
the recent literature and contributes a method together with an evaluation protocol.

\section{Related Work (real literature)}
The following works, retrieved from public bibliographic APIs, frame our study:
\begin{itemize}
\item Sambrook et al. (1989) — "Molecular Cloning: A Laboratory Manual" (cited 130003).
\item Network (2012) — "Comprehensive molecular portraits of human breast tumours", Nature (cited 12476).
\item McLendon et al. (2008) — "Comprehensive genomic characterization defines human glioblastoma genes and core pathways", Nature (cited 7831).
\item Horvath (2013) — "DNA methylation age of human tissues and cell types", Genome biology (cited 7417).
\item Hegi et al. (2005) — "<i>MGMT</i> Gene Silencing and Benefit from Temozolomide in Glioblastoma", New England Journal of Medicine (cited 7299).
\item Bird (2002) — "DNA methylation patterns and epigenetic memory", Genes \& Development (cited 7144).
\end{itemize}

\section{Method}
We model distribution shift explicitly and optimise a worst-case objective over an uncertainty set of plausible test distributions. We instantiate this idea for Epigenetic Dna Methylation and analyse its cost/quality trade-off.

\section{Experiments (Illustrative)}
\begin{center}
\begin{tabular}{lcc}
System & Primary Metric & Efficiency \\ \hline
Strong baseline & 0.812 & 1.00$\times$ \\
Ablation & 0.828 & 0.92$\times$ \\
\textbf{Ours} & \textbf{0.851} & \textbf{0.71$\times$}
\end{tabular}
\end{center}
\emph{Metrics simulated to illustrate the intended effect; not measured.}

\section{Conclusion}
We connected a real literature to a concrete method for Epigenetic Dna Methylation. Future work will
validate the illustrative results with real experiments.

\begin{thebibliography}{9}
\bibitem{ref1} Joseph Sambrook, Elisabeth Fritsch. Molecular Cloning: A Laboratory Manual. \emph{}, 1989. 
\bibitem{ref2} The Cancer Genome Atlas Network. Comprehensive molecular portraits of human breast tumours. \emph{Nature}, 2012. DOI:10.1038/nature11412
\bibitem{ref3} Roger E. McLendon, D D Bigner, Allan H. Friedman et al.. Comprehensive genomic characterization defines human glioblastoma genes and core pathways. \emph{Nature}, 2008. DOI:10.1038/nature07385
\bibitem{ref4} Steve Horvath. DNA methylation age of human tissues and cell types. \emph{Genome biology}, 2013. DOI:10.1186/gb-2013-14-10-r115
\bibitem{ref5} Monika E. Hegi, Annie‐Claire Diserens, Thierry Gorlia et al.. <i>MGMT</i> Gene Silencing and Benefit from Temozolomide in Glioblastoma. \emph{New England Journal of Medicine}, 2005. DOI:10.1056/nejmoa043331
\bibitem{ref6} Adrian Bird. DNA methylation patterns and epigenetic memory. \emph{Genes \& Development}, 2002. DOI:10.1101/gad.947102
\bibitem{ref7} Rudolf Jaenisch, Adrian Bird. Epigenetic regulation of gene expression: how the genome integrates intrinsic and environmental signals. \emph{Nature Genetics}, 2003. DOI:10.1038/ng1089
\bibitem{ref8} Peter A. Jones. Functions of DNA methylation: islands, start sites, gene bodies and beyond. \emph{Nature Reviews Genetics}, 2012. DOI:10.1038/nrg3230
\end{thebibliography}
\end{document}
